% DOTE2011G · In-Class Exercise 1 · Instructor answer key
% Do not distribute to students before the deadline.
% Compile with: xelatex Ex1-solutions.tex

\documentclass[11pt,a4paper,UTF8]{ctexart}

\usepackage[margin=2.2cm]{geometry}
\usepackage{amsmath,amssymb}
\usepackage{enumitem}
\usepackage{xcolor}
\usepackage{tcolorbox}

\definecolor{lantern}{HTML}{C43C00}
\definecolor{night}{HTML}{0B1C3A}

\begin{document}

\begin{center}
  {\Large\bfseries\color{night} In-Class Exercise 1 · Solutions}\\[0.3em]
  {\normalsize Full house counting (instructor key --- not for students)}
\end{center}

\vspace{0.6em}

\noindent\textbf{Riddle meaning.}
A full house: three cards of one rank and two cards of a different rank.

\vspace{0.5em}

\noindent\textbf{Count.}
\[
\begin{aligned}
N
&= C(13,1)\,C(4,3)\,C(12,1)\,C(4,2) \\
&= 13 \times 4 \times 12 \times 6 \\
&= 3744.
\end{aligned}
\]

\begin{itemize}[leftmargin=1.4em]
  \item $C(13,1)$: choose the rank for the three-of-a-kind.
  \item $C(4,3)$: choose 3 suits out of 4 for that rank.
  \item $C(12,1)$: choose a different rank for the pair.
  \item $C(4,2)$: choose 2 suits out of 4 for the pair.
\end{itemize}

\noindent\textit{Note.}
Do not multiply by an extra factor for ``which block is the triple / which is the pair'':
the two ranks are already chosen in order (triple rank first, pair rank second), so each full house is counted once.

\vspace{0.55em}

\noindent\textbf{Optional probability.}
\[
C(52,5)=2\,598\,960,
\qquad
P(\text{full house})=\frac{3744}{C(52,5)}=\frac{3744}{2\,598\,960}=\frac{6}{4165}.
\]
(The simplified fraction is optional for students.)

\vspace{0.7em}
\begin{tcolorbox}[colback=lantern!5,colframe=lantern,title=\textbf{Quick check}]
Final answer: $\mathbf{3744}$ full-house hands.
\end{tcolorbox}

\end{document}
